% ==============================================================================
% Fichier     : chapitres/02_methodo_juridique.tex
% Titre       : Cadre Juridique et Méthodologies de Tests d'Intrusion
% Auteur      : MINKA MI NGUIDJOÏ Thierry Emmanuel
% ==============================================================================

\chapter{Cadre Juridique et Méthodologies de Tests d'Intrusion}
\label{chap:02_cadre_juridique}

Le Hacking Éthique n'est pas une absence de règles, mais une pratique strictement
encadrée. Ce chapitre détaille les limites légales à ne jamais franchir, tant dans
le droit français que dans le droit camerounais, et les standards méthodologiques
à suivre pour garantir la qualité et la valeur probatoire d'une mission.

% ==============================================================================
\section{Le Cadre Juridique}
\label{sec:cadre_juridique}
% ==============================================================================

\subsection{Droit Français : La Loi Godfrain}
\label{subsec:loi_godfrain}

En France, le cadre légal principal est défini par le \textbf{Code pénal}
(Articles 323-1 à 323-7), communément appelé \termecle{Loi Godfrain} (du nom
du parlementaire qui en fut le rapporteur, 1988).

\begin{boxalert}
\textbf{Régime de sanctions, Code pénal français :}
\begin{itemize}
    \item \textbf{Accès ou maintien frauduleux} (Art.~323-1) : jusqu'à
          2~ans d'emprisonnement et 30\,000~€ d'amende ;
          jusqu'à 3~ans et 45\,000~€ si cela a entraîné modification
          ou suppression de données.
    \item \textbf{Entrave au fonctionnement} (Art.~323-2) : jusqu'à
          5~ans et 150\,000~€ d'amende.
    \item \textbf{Introduction frauduleuse de données} (Art.~323-3) :
          jusqu'à 5~ans et 150\,000~€.
    \item \textbf{Circonstances aggravantes} (Art.~323-4 à 323-7) :
          peines doublées en cas d'organisation criminelle ou d'atteinte
          à un \sigle{SIIV} (Système d'Information d'Importance Vitale).
\end{itemize}
\textit{Sans autorisation écrite préalable, toute intrusion, même
sans intention malveillante, constitue un délit.}
\end{boxalert}

\subsection{Droit Camerounais : Loi sur la Cybersécurité}
\label{subsec:loi_cameroun}

Dans le contexte africain francophone, et en particulier au Cameroun, le cadre
légal de référence est défini par deux textes complémentaires.

\begin{boxalert}
\textbf{Régime de sanctions, Droit camerounais :}

\medskip
\textbf{Loi n\textsuperscript{o}~2010/012 du 21~décembre~2010} relative à la
cybersécurité et à la cybercriminalité :
\begin{itemize}
    \item \textbf{Accès non autorisé à un système} (Art.~72) : jusqu'à
          2~ans d'emprisonnement et 5~000\,000~FCFA d'amende.
    \item \textbf{Atteinte à l'intégrité d'un système} (Art.~73) :
          jusqu'à 5~ans et 10\,000\,000~FCFA.
    \item \textbf{Interception illégale de données} (Art.~74) :
          jusqu'à 2~ans et 5\,000\,000~FCFA.
    \item \textbf{Circonstances aggravantes} : peines doublées lorsque
          les faits sont commis contre des infrastructures d'État ou des
          systèmes financiers.
\end{itemize}

\medskip
\textbf{Loi n\textsuperscript{o}~2010/013 du 21~décembre~2010} sur les
communications électroniques : encadre les obligations des opérateurs
et prestataires de services numériques ; toute prestation de sécurité
doit être déclarée à l'\sigle{ANTIC}.

\medskip
\textbf{Autorité de régulation :} l'\sigle{ANTIC} (Agence Nationale des
Technologies de l'Information et de la Communication) est l'organisme
compétent en matière de cybersécurité au Cameroun. Toute mission de
pentest réalisée dans un cadre professionnel au Cameroun doit, selon les
cas, faire l'objet d'une déclaration préalable.
\end{boxalert}

\subsection{Tableau Comparatif France / Cameroun}
\label{subsec:comparatif}

\begin{table}[htbp]
\centering
\caption{Comparatif des régimes juridiques applicables au pentest}
\label{tab:comparatif_juridique}
\begin{tabular}{p{3.8cm}p{4.8cm}p{4.8cm}}
\toprule
\textbf{Critère}
    & \textbf{France (Loi Godfrain)}
    & \textbf{Cameroun (Loi 2010/012)} \\
\midrule
Texte de référence
    & Code pénal, Art.~323-1 à 323-7
    & Loi n\textsuperscript{o}~2010/012 \\
\addlinespace
Accès non autorisé
    & 2~ans / 30\,000~€
    & 2~ans / 5\,000\,000~FCFA \\
\addlinespace
Atteinte à l'intégrité
    & 5~ans / 150\,000~€
    & 5~ans / 10\,000\,000~FCFA \\
\addlinespace
Autorité de régulation
    & ANSSI
    & ANTIC \\
\addlinespace
Déclaration obligatoire
    & Non (sauf \sigle{SIIV})
    & Selon nature de la prestation \\
\addlinespace
Condition de légalité
    & Autorisation écrite + périmètre
    & Autorisation écrite + déclaration ANTIC \\
\bottomrule
\end{tabular}
\end{table}

\begin{boxinfo}
\textbf{Application pratique pour le cours :}
Les travaux pratiques de ce cours sont réalisés sur un environnement de
laboratoire totalement isolé. Ils ne sont soumis à aucune des obligations
déclaratives mentionnées ci-dessus. En revanche, tout étudiant réalisant
une mission de pentest professionnelle au Cameroun est tenu de connaître
et respecter le cadre de la loi 2010/012 et les exigences de l'\sigle{ANTIC}.
\end{boxinfo}

% ==============================================================================
\section{Les Trois Conditions Cumulatives de Légalité}
\label{sec:conditions_legalite}
% ==============================================================================

Quelle que soit la juridiction applicable, la légalité d'un test d'intrusion
repose sur trois conditions \textbf{cumulatives}, l'absence de l'une suffit à
caractériser l'infraction.

\begin{enumerate}
    \item \textbf{Autorisation expresse et écrite :} Un document contractuel
          (lettre de mission, bon de commande, avenant) signé par le
          \textbf{propriétaire légitime} du système ou son mandataire dûment
          habilité. Une autorisation verbale n'a aucune valeur juridique.

    \item \textbf{Définition précise du périmètre :} L'autorisation doit
          spécifier exactement quelles machines, plages \sigle{IP}, applications
          ou services sont testés. La formulation \og{}tester tout le réseau\fg{}
          est trop vague et ne protège pas le pentester : si une machine
          hors-périmètre est atteinte par rebond, l'infraction est constituée.

    \item \textbf{Respect des Règles d'Engagement :} Respecter scrupuleusement
          les horaires, les méthodes interdites (DoS, destruction de données)
          et les cibles exclues définies dans le document \sigle{ROE} signé
          (voir Section~\ref{sec:roe}).
\end{enumerate}

\begin{boxinfo}
\textbf{Cas particulier, Bug Bounty :} Les programmes de \textit{bug bounty}
(HackerOne, Bugcrowd) constituent une forme d'autorisation implicite encadrée
par les \textit{Terms of Service} du programme. Ils ne dispensent pas du
respect des conditions 2 et 3 : le périmètre (\textit{in-scope assets}) et
les règles d'engagement sont définis dans le programme et s'imposent au
chercheur.
\end{boxinfo}

% ==============================================================================
\section{Déontologie du Hacker Éthique}
\label{sec:deontologie}
% ==============================================================================

L'éthique professionnelle est ce qui distingue le pentester du cybercriminel.
Elle ne se réduit pas à respecter la loi, elle impose des obligations positives
vis-à-vis du client, des tiers et de la profession.

\begin{itemize}
    \item \textbf{Confidentialité :} Les données collectées lors de la mission
          (mots de passe, données personnelles, stratégie commerciale, secrets
          industriels) ne doivent \textbf{jamais} être divulguées, conservées
          après la mission ou utilisées à des fins personnelles. L'accord de
          confidentialité (\sigle{NDA}) est obligatoire.

    \item \textbf{Intégrité :} Ne pas détruire de données, ne pas modifier les
          configurations au-delà de ce qui est nécessaire à la démonstration,
          ne pas installer de \textit{backdoors} persistantes sans accord
          contractuel explicite.

    \item \textbf{Divulgation responsable (\textit{Responsible Disclosure}) :}
          Si une vulnérabilité critique est découverte sur un système hors
          périmètre (par rebond), le pentester a l'obligation déontologique —
          et parfois légale, d'en informer le propriétaire concerné, via
          un processus de divulgation coordonnée, sans l'exploiter.

    \item \textbf{Non-nuisance :} Le pentester doit évaluer le risque de chaque
          action avant de l'exécuter. Une exploitation qui risque de rendre
          un service indisponible doit être discutée avec le client avant
          d'être tentée, même si elle est techniquement \textit{in-scope}.

    \item \textbf{Objectivité du rapport :} L'objectif est d'aider le client
          à se sécuriser, non de démontrer sa propre expertise technique. Un
          rapport qui gonfle la criticité des vulnérabilités pour impressionner
          est contraire à l'éthique professionnelle.

    \item \textbf{Formation continue :} Le professionnel est responsable de
          maintenir ses compétences à jour. Utiliser des techniques obsolètes
          ou passer à côté de vecteurs d'attaque actuels constitue un manque
          à son obligation de moyens vis-à-vis du client.
\end{itemize}

% ==============================================================================
\section{Méthodologies Standardisées}
\label{sec:methodologies}
% ==============================================================================

Ne pas utiliser de méthodologie formalisée, c'est risquer d'oublier des failles
critiques et de produire un rapport non défendable devant un client exigeant.
Cinq standards structurent aujourd'hui l'industrie du pentest.

\subsection{PTES, \textit{Penetration Testing Execution Standard}}
\label{subsec:ptes}

\sigle{PTES} est la méthodologie la plus utilisée pour les tests techniques.
Elle divise le pentest en \textbf{7~phases} et constitue le fil conducteur de
ce cours.

\begin{table}[htbp]
\centering
\caption{Phases PTES avec mapping MITRE ATT\&CK et chapitres du cours}
\label{tab:ptes_mitre}
\begin{tabular}{p{0.4cm}p{3.5cm}p{3.5cm}p{5cm}}
\toprule
\textbf{\#} & \textbf{Phase PTES} & \textbf{MITRE ATT\&CK} & \textbf{Chapitre du cours} \\
\midrule
1 & \textit{Pre-Engagement}
  & N/A
  & Chap.~\ref{chap:mandat} (ROE, scope) \\
\addlinespace
2 & \textit{Intelligence Gathering}
  & TA0043 Reconnaissance
  & Chap.~\ref{chap:osint} (OSINT) \\
\addlinespace
3 & \textit{Threat Modeling}
  & TA0043, TA0042
  & Chap.~\ref{chap:osint} (modélisation) \\
\addlinespace
4 & \textit{Vulnerability Analysis}
  & TA0043
  & Chap.~\ref{chap:scanning} (scanning) \\
\addlinespace
5 & \textit{Exploitation}
  & TA0001, TA0002
  & Chap.~\ref{chap:exploit} \\
\addlinespace
6 & \textit{Post-Exploitation}
  & TA0003--TA0010
  & Chap.~\ref{chap:postexploit}, \ref{chap:evasion} \\
\addlinespace
7 & \textit{Reporting}
  & N/A
  & Chap.~\ref{chap:reporting} \\
\bottomrule
\end{tabular}
\end{table}

\subsection{OSSTMM, \textit{Open Source Security Testing Methodology Manual}}
\label{subsec:osstmm}

Plus théorique et exhaustive que \sigle{PTES}, l'\sigle{OSSTMM} se concentre
sur la \textbf{rigueur scientifique} de l'audit et la mesure quantitative de
la surface d'attaque (\textit{Attack Surface}). Elle introduit le concept de
\termecle{RAV} (\textit{Risk Assessment Value}), permettant de calculer un
score de sécurité objectif et reproductible. Elle couvre également les vecteurs
non-informatiques : sécurité physique, ingénierie sociale, sécurité des
télécommunications. Elle est souvent privilégiée pour les audits complets
d'une organisation.

\subsection{NIST SP 800-115}
\label{subsec:nist}

Guide technique du \sigle{NIST} (\textit{National Institute of Standards and
Technology}) américain. Très structuré, il est souvent contractuellement exigé
pour les organismes publics américains ou les prestataires soumis aux exigences
\sigle{FedRAMP}/\sigle{FISMA}. Sa Section~3 sur la planification et la Section~7
sur la production du rapport sont applicables universellement et constituent
d'excellentes références.

\subsection{OWASP, \textit{Open Web Application Security Project}}
\label{subsec:owasp}

L'\sigle{OWASP} publie deux documents complémentaires qu'il faut impérativement
distinguer :

\begin{itemize}
    \item \textbf{OWASP Top~10} : liste des dix catégories de vulnérabilités
          Web les plus critiques, mise à jour périodiquement (dernière édition :
          2021). C'est un \textbf{référentiel de sensibilisation}, non un guide
          opérationnel. Il sera utilisé au Chapitre~\ref{chap:exploit}.

    \item \textbf{OWASP Testing Guide v4.2} (\sigle{OTG}) : guide opérationnel
          de plus de 500~pages détaillant les procédures de test pour chaque
          catégorie de vulnérabilité. C'est la référence technique pour les
          tests d'applications web, il précise \textit{comment} tester, pas
          seulement \textit{quoi} chercher. Indispensable pour les missions
          de pentest applicatif.
\end{itemize}

\subsection{MITRE ATT\&CK}
\label{subsec:mitre}

Bien que \sigle{MITRE ATT\&CK} ne soit pas une méthodologie de pentest à
proprement parler, il s'est imposé comme le \textbf{langage universel} de
description des techniques d'attaque. Son utilisation dans ce cours est
systématique : chaque technique présentée est accompagnée de son identifiant
\sigle{ATT\&CK} (ex.~: T1190, T1558.003), facilitant la communication avec
les équipes de défense et la production de rapports compréhensibles par les
équipes \sigle{SOC}.

\begin{boxinfo}
\textbf{Structure MITRE ATT\&CK :} Le framework est organisé en
\textbf{14~tactiques} (le \textit{pourquoi} de l'action), chacune
décomposée en \textbf{techniques} (le \textit{comment}) et
\textbf{sous-techniques} (les variantes précises). Exemple :
\textit{Tactique}~TA0006 (Credential Access) →
\textit{Technique}~T1003 (OS Credential Dumping) →
\textit{Sous-technique}~T1003.001 (LSASS Memory), ce qui correspond
à l'utilisation de Mimikatz au Chapitre~\ref{chap:postexploit}.
\end{boxinfo}

% ==============================================================================
\section{Types de Tests d'Intrusion}
\label{sec:types_tests}
% ==============================================================================

Le niveau d'information fourni au pentester avant la mission définit le
\termecle{type de test} et la nature de la menace simulée.

\begin{table}[htbp]
\centering
\caption{Types de tests d'intrusion}
\label{tab:types_tests}
\begin{tabular}{p{2.2cm}p{3.8cm}p{3.8cm}p{3cm}}
\toprule
\textbf{Type}
    & \textbf{Information fournie}
    & \textbf{Menace simulée}
    & \textbf{Usage typique} \\
\midrule
\textbf{Black Box}
    & Aucune (0~knowledge), seule la cible est connue
    & Attaquant externe inconnu, cybercriminel opportuniste
    & Test périmétrique, audit externe \\
\addlinespace
\textbf{Gray Box}
    & Partielle, compte utilisateur standard, plages IP
    & Insider malveillant, partenaire compromis, attaquant post-phishing
    & Audit interne, simulation réaliste \\
\addlinespace
\textbf{White Box}
    & Totale, code source, schémas réseau, comptes admin
    & Auditeur interne, développeur malveillant
    & Audit de code, revue d'architecture \\
\addlinespace
\textbf{Crystal Box}
    & Totale + accès collaboratif aux équipes internes
    & N/A (mode coopératif)
    & Red Team / Blue Team conjoint \\
\bottomrule
\end{tabular}
\end{table}

\begin{boxlizbiz}
\textbf{Choix pour LiZbiZ :}
La direction a opté pour un test \textbf{Gray Box}. Vous disposez d'un compte
utilisateur standard (\og{}stagiaire\fg{}) sur l'intranet, simulant un employé
peu qualifié ou un attaquant ayant réussi un phishing initial. Cette configuration
est la plus représentative des scénarios réels : selon les statistiques \sigle{DBIR}
(Verizon Data Breach Investigations Report), plus de 74\% des compromissions
impliquent un facteur humain en phase initiale.
\end{boxlizbiz}

% ==============================================================================
\section{Le Cycle de Vie d'une Attaque}
\label{sec:cycle_vie}
% ==============================================================================

Indépendamment de la méthodologie choisie, le cycle de vie d'une intrusion suit
toujours un flux logique. La compréhension de ce cycle est fondamentale : elle
permet au pentester d'anticiper les actions suivantes et au défenseur de placer
ses contrôles aux points de rupture les plus efficaces.

\begin{figure}[htbp]
\centering
\begin{tikzpicture}[
    node distance = 1.5cm and 2.4cm,
    phase/.style  = {
        rectangle, draw, thick, rounded corners=6pt,
        minimum width=2.8cm, minimum height=1cm,
        align=center, font=\small\bfseries
    },
    ta/.style = {
        font=\tiny, text=gray, align=center
    }
]
    % Nœuds, deux rangées
    \node[phase, fill=blue!10]   (recon)    {Reconnaissance\\{\tiny TA0043}};
    \node[phase, fill=blue!15,   right=of recon]  (scan)     {Scanning\\{\tiny TA0043}};
    \node[phase, fill=orange!20, right=of scan]   (exploit)  {Exploitation\\{\tiny TA0001-02}};
    \node[phase, fill=red!15,    below=of exploit] (postex)   {Post-Exploitation\\{\tiny TA0003-08}};
    \node[phase, fill=red!10,    left=of postex]  (maintain) {Maintien d'Accès\\{\tiny TA0003}};
    \node[phase, fill=gray!20,   left=of maintain] (cover)    {Couverture\\des Traces\\{\tiny TA0005}};

    % Flèches principales
    \draw[->, thick] (recon)   -- (scan);
    \draw[->, thick] (scan)    -- (exploit);
    \draw[->, thick] (exploit) -- (postex);
    \draw[->, thick] (postex)  -- (maintain);
    \draw[->, thick] (maintain)-- (cover);

    % Boucle : nouvelle info remonte à la reconnaissance
    \draw[->, thick, dashed, CouleurAccent]
        (maintain.south)
        to[bend right=45]
        node[below, font=\footnotesize\itshape] {nouvelle info}
        (recon.south);
\end{tikzpicture}
\caption{Cycle de vie d'une attaque, phases et tactiques MITRE ATT\&CK associées}
\label{fig:cycle_vie}
\end{figure}

\begin{table}[htbp]
\centering
\caption{Correspondance phases du cycle de vie et chapitres du cours}
\label{tab:cycle_chapitres}
\begin{tabular}{p{3.2cm}p{3cm}p{6.5cm}}
\toprule
\textbf{Phase} & \textbf{MITRE ATT\&CK} & \textbf{Outils / Techniques clés} \\
\midrule
Reconnaissance
    & TA0043
    & Whois, dig, Shodan, Google Dorks, theHarvester \\
\addlinespace
Scanning
    & TA0043
    & Nmap, Dirb, Nikto, Nessus/OpenVAS \\
\addlinespace
Exploitation
    & TA0001, TA0002
    & SQLMap, Metasploit, Hydra, EternalBlue \\
\addlinespace
Post-Exploitation
    & TA0003--TA0008
    & Mimikatz, BloodHound, CrackMapExec, Chisel \\
\addlinespace
Maintien d'Accès
    & TA0003
    & \textit{Web Shell}, tâches planifiées, persistance AD \\
\addlinespace
Couverture des Traces
    & TA0005
    & Wevtutil, \textit{timestomping}, techniques \textit{fileless} \\
\bottomrule
\end{tabular}
\end{table}

\begin{boxlizbiz}
\textbf{Illustration LiZbiZ, lecture du cycle :}
\begin{enumerate}
    \item \textbf{Reconnaissance} : collecte OSINT révèle
          \texttt{admin@lizbiz.com} et un Apache~2.2.22 non patché.
    \item \textbf{Scanning} : Nmap confirme le port~80 ouvert ; Dirb découvre
          \texttt{/backup/old\_passwd.txt}.
    \item \textbf{Exploitation} : injection SQL sur le formulaire de connexion
          + exploitation \texttt{MS17-010} sur le poste Windows interne.
    \item \textbf{Post-Exploitation} : escalade via \texttt{sudo vim},
          dump Mimikatz, Pass-the-Hash vers le \sigle{DC}.
    \item \textbf{Maintien d'Accès} : dépôt d'un \textit{web shell} \sigle{PHP}
          sur le serveur \sigle{DMZ} (simulé uniquement).
    \item \textbf{Couverture} : effacement des logs documenté dans le rapport
          (non exécuté en conditions réelles sans accord client).
\end{enumerate}
\end{boxlizbiz}

% ==============================================================================
\section*{Points Clés à Retenir}
% ==============================================================================

\begin{boxinfo}
\textbf{Ce qu'il faut retenir de ce chapitre :}
\begin{enumerate}
    \item La légalité d'un pentest repose sur \textbf{trois conditions
          cumulatives} : autorisation écrite, périmètre précis, respect
          des ROE.
    \item Le cadre camerounais (loi 2010/012) impose des sanctions comparables
          au droit français et exige une \textbf{connaissance de l'\sigle{ANTIC}}
          pour les missions professionnelles.
    \item \textbf{PTES} structure les 7~phases ; \textbf{MITRE ATT\&CK} en
          fournit le vocabulaire universel.
    \item \textbf{OWASP Top~10} ≠ \textbf{OWASP Testing Guide} : le premier
          liste les risques, le second explique comment les tester.
    \item Le type de test (\textit{Black/Gray/White Box}) détermine la menace
          simulée et le périmètre de connaissance du pentester.
    \item La déontologie (\sigle{NDA}, divulgation responsable,
          non-nuisance) est aussi contraignante que la loi.
\end{enumerate}
\end{boxinfo}

% ==============================================================================
\section*{Exercice Pratique}
% ==============================================================================

\begin{boxexo}
\textbf{Exercice 2.1, Analyse d'un cas juridique}

\medskip
\textbf{Scénario :} Un consultant en sécurité reçoit un email de la part d'un
\sigle{DSI} lui demandant verbalement de tester la sécurité du site web de son
entreprise. Il procède au scan et découvre une injection SQL critique.
Cependant, le scan a également atteint un sous-domaine hébergé chez un
prestataire tiers.

\medskip
\textbf{Questions :}
\begin{enumerate}
    \item Quelle(s) condition(s) de légalité manquent dans ce scénario ?
    \item Quelles infractions le consultant risque-t-il au regard du Code pénal
          français ? Et de la loi camerounaise 2010/012 ?
    \item Que doit faire le consultant concernant la vulnérabilité découverte
          sur le sous-domaine tiers ?
    \item Rédigez en 10~lignes la clause d'autorisation minimale qui aurait
          permis d'éviter ce problème.
\end{enumerate}

\medskip
\textbf{Livrable :} Réponses argumentées + clause rédigée (à déposer
avant le début du TP Chapitre~\ref{chap:osint}).
\end{boxexo}

\begin{boxexo}
\textbf{Exercice 2.2, Mapping MITRE ATT\&CK}

\medskip
Pour chacune des techniques suivantes, identifier la tactique MITRE ATT\&CK
parente, l'identifiant de technique et de sous-technique (si applicable), et
le chapitre du cours où elle est traitée :
\begin{enumerate}
    \item Extraction de mots de passe depuis la mémoire \sigle{LSASS}.
    \item Exploitation d'une injection SQL pour accéder à la base de données.
    \item Utilisation de \texttt{certutil.exe} pour télécharger un fichier.
    \item Kerberoasting sur un compte de service Active Directory.
    \item Effacement du journal de sécurité Windows.
\end{enumerate}

\medskip
\textbf{Ressource :} \texttt{https://attack.mitre.org}
\end{boxexo}
